\section{Design of the empirical network}\label{app:database-design}

\subsection{What the reconstruction needs to measure}

The empirical object is a matrix of debt claims linking sectors and countries at the end of 2024. Its purpose is to identify who loses receipts when a sovereign or a bank repays less, and how those losses affect other banks' ability to pay. This requires consistent creditor and debtor positions, sovereign holdings by banks, and a reconciliation with bank liabilities and equity. It does not require a microeconomic model of how each portfolio was chosen or a complete network of every financial instrument.

No public source supplies this matrix directly. Banking statistics provide strong controls on resident bank balance sheets but incomplete bilateral detail. Financial accounts cover other sectors but often aggregate foreign counterparties. Portfolio surveys provide cross-border securities positions with uneven sector coverage. The construction therefore combines these sources in a hierarchy: published instrument totals anchor the accounts, available bilateral observations guide their allocation, and missing detail is filled within the applicable controls. The exposure register distinguishes these cases and retains the adjustments.

\subsection{The main choices, especially for banks}

\paragraph{A common residence perimeter.}
Bank nodes represent offices located in each country, including foreign-owned offices and excluding the foreign affiliates of domestic groups. This matches the residence basis of the financial accounts and public securities positions. Bank for International Settlements (BIS) consolidated claims instead describe worldwide banking groups; combining them directly with resident funding and equity would mix institutions and foreign affiliates across the two sides of the balance sheet. We consequently use European Central Bank (ECB) resident-bank statistics as the bank controls. This preserves a common statistical perimeter without claiming that every bilateral position is observed.

\paragraph{Bank assets and funding from their own source controls.}
Public bank loan data distinguish firms and households, subject to the reporting-sector adjustments described below. Household bank funding comes from depositor-sector observations and securities holdings. Equity follows the published capital accounts, with a matched asset adjustment for provisions. These stocks are measured separately: missing funding classifications remain explicit, and the exposure matrix is subsequently reconciled to the complete balance sheet. This keeps creditor incidence and bank loss absorption tied to their relevant public sources.

\paragraph{Sovereign exposures anchored in bank portfolios.}
Aggregate bank government-security holdings are controlled by ECB balance-sheet statistics. Domestic securities holdings at face value are preferred where public; Italy requires an allocation within its bank portfolio control. Foreign holdings use residence-based portfolio information and bilateral securities envelopes. These allocations must respect both bank instrument controls and issuer totals. Public government debt-security stocks anchor the sovereign columns, while valuation differences and adjustments remain recorded. Matching the total debt stock does not make every bilateral bank holding an independently observed face amount.

\paragraph{Remove identified intragroup claims and retain other domestic links.}
Related-office loan margins permit removal of identified intragroup positions, within the reconstructed gross exposures and with matched bank balance-sheet adjustments. The remaining domestic claims are retained as gross obligations among distinct banks. The result is neither a worldwide consolidated group network nor a fully consolidated national banking sector. Public data do not identify every intragroup security position or foreign-office link. The model's common repayment rate within each country is a separate aggregation assumption, whose sensitivity is reported.

\paragraph{Incomplete links remain distinguishable from observations.}
Missing bilateral detail is allocated using published margins, available upper bounds and stated priors. Unidentified holders and issuers have explicit reporting accounts. Loan assets and their funding mirrors are reconciled rather than added as separate claims; bank securities drawn from different sources are likewise not counted twice. Assets outside the selected debt instruments and unresolved coverage or valuation differences enter the bank balance-sheet reconciliation. They are not treated as independently measured bilateral exposures.

The resulting database is a controlled reconstruction, not an observation of the complete underlying network. Its strength is a consistent public-data account of the positions relevant to repayment, with imputations and limits visible. Its principal uncertainties concern bilateral allocations, valuation and the treatment of intragroup and gross domestic claims. Prudential headroom, private-credit crisis losses and default costs enter only in the subsequent model mapping; they do not determine the matrix entries.

\subsection{How to read the construction}

Section~\ref{app:database_construction} gives the operational sequence: define coverage and source controls; construct bank loans and funding; reconcile interbank claims and related offices; allocate securities holdings; and check issuer totals and unknown counterparties. Section~\ref{app:residence-model-mapping} then reconciles these selected claims with full bank balance sheets and explains the crisis specification. This separates what is measured or imputed from how a given loss changes repayment.
